% =========================================================
%  sesi/04/sesi.tex
%  Sesi 4 — Audio dan Trigger Interaksi
% =========================================================

\sesi{04}{Audio dan Trigger Interaksi}{4.2}{130 menit}

Ini pertemuan minggu keempat belas. Level Rak sudah tampak lengkap: gambar
pemain berganti menurut geraknya, skor bertambah, dan nyawa berkurang.
Namun seluruhnya masih bisu. Pemain yang melompat tidak mendengar apa pun,
dan artefak yang lenyap tidak memberi tanggapan selain angka pada sudut layar.

Sesi ini menambahkan jalan umpan balik yang kedua, yaitu bunyi. Berkas bunyi
diimpor dan dipisahkan menjadi musik latar dan efek suara, peristiwa permainan
dihubungkan dengan bunyinya melalui satu skrip pemusat, lalu waktu pemanggilan
dan tingkat kekuatan disesuaikan. Partikel sederhana ditambahkan sebagai
pelengkap bila waktu memungkinkan. Mekanik, animasi, dan antarmuka dari sesi 3
tidak diubah; setiap kali sebuah bunyi dipasang, jalankan kembali scene dan
pastikan pemain masih dapat berjalan, melompat, mengambil artefak, dan membuka
pintu seperti sebelumnya.

% ---------------------------------------------------------
\bagian{Tujuan Sesi}

\begin{tujuanbox}
Setelah menuntaskan sesi ini, Anda dapat menunjukkan:
\begin{enumerate}[leftmargin=1.5em, itemsep=0.2em]
  \item berkas \istilah{AudioClip} yang terimpor dengan pengaturan
        \menu{Load Type} dan \menu{Compression} yang tercatat;
  \item objek \code{ManajerBunyi} berisi dua \istilah{AudioSource} terpisah
        untuk \istilah{BGM} dan \istilah{SFX};
  \item bunyi lompat, bunyi mengambil artefak, dan bunyi kalah yang terdengar
        pada peristiwanya masing-masing;
  \item \istilah{AudioMixer} dengan tingkat \code{Musik} dan \code{Efek} yang
        kekuatannya dapat diatur terpisah;
  \item catatan percobaan volume dan \istilah{timing} beserta alasan setiap
        nilai yang dipilih.
\end{enumerate}
\end{tujuanbox}

% ---------------------------------------------------------
\bagian{Prasyarat}

\begin{prasyaratbox}
\begin{itemize}[leftmargin=1.4em, itemsep=0.25em]
  \item Checkpoint sesi 3 sudah berjalan: pemain beranimasi, skor dan nyawa
        tampil pada \istilah{Canvas}, dan pintu menolak pemain selama artefak
        belum lengkap.
  \item Anda telah membaca Bab XV bagian 15.1 sampai 15.5, yaitu bagian
        tentang peran bunyi, \istilah{AudioClip} dan \istilah{AudioSource},
        peristiwa bunyi, pengaturan waktu dan kekuatan, serta partikel.
  \item Anda memiliki satu berkas musik latar dan tiga berkas efek suara
        pendek: lompat, mengambil benda, dan kalah. Catat sumber serta
        lisensinya sejak awal.
  \item Scene sesi 3 tersimpan. Sesi ini menimpa scene itu, sehingga simpanlah
        salinan bila Anda masih memerlukannya.
\end{itemize}
\end{prasyaratbox}

% ---------------------------------------------------------
\bagian{Peta Kemampuan}

\petasesi
  {Level yang tampak dan beranimasi, tetapi tidak mengeluarkan bunyi sama
   sekali.}
  {\istilah{AudioClip}, \istilah{AudioSource}, \istilah{BGM},
   \istilah{SFX}, \istilah{AudioMixer}, dan pemanggilan bunyi pada peristiwa.}
  {Pengujian fungsional, \istilah{playtest} antar kelompok, perbaikan
   \istilah{bug}, dan demo sesi 5.}

% ---------------------------------------------------------
\bagian{Konsep yang Dipakai}

\subsection*{Berkas dan pemutarnya adalah dua hal}

\istilah{AudioClip} adalah berkas bunyi yang telah diimpor. Ia hanya berisi
data; ia tidak tahu di mana ia berada dan tidak dapat berbunyi sendiri.
\istilah{AudioSource} adalah komponen yang menempel pada \istilah{GameObject}
dan memutar klip itu. Pertanyaan ``di mana bunyi ini terdengar'' selalu
dijawab oleh letak \code{AudioSource}, bukan oleh berkasnya.

Karena itu, bunyi yang berbeda memerlukan sumber yang berbeda pula. Musik
latar dan efek suara tidak boleh memakai satu \code{AudioSource}, sebab setiap
efek suara akan menghentikan musik yang sedang berjalan. Sesi ini memakai dua
sumber: \code{musicSource} untuk musik dan \code{sfxSource} untuk efek.

\subsection*{Musik berulang, efek sekali pakai}

\istilah{BGM} berjalan terus dan diputar berulang, sehingga ia dipanggil
dengan \code{Play()} setelah \code{loop} dinyalakan. \istilah{SFX} dipanggil
menurut kejadian dan harus dapat bertumpuk; dua artefak yang diambil berdekatan
menghasilkan dua bunyi yang hampir bersamaan. Untuk itu dipakai
\code{PlayOneShot()}, yang tidak memotong bunyi sebelumnya. Memakai
\code{Play()} untuk efek suara menghasilkan bunyi yang terputus setiap kali
panggilan baru datang.

\subsection*{Bunyi dipanggil pada peristiwa, bukan pada keadaan}

Bunyi lompat dipanggil ketika tombol ditekan, bukan selama tombol ditahan.
Pemeriksaannya dilakukan dengan \code{Input.GetButtonDown} di \code{Update},
tempat peristiwa masukan dibaca. Bunyi artefak dipanggil dari
\istilah{trigger} yang sudah ada sejak sesi 2, dan bunyi kalah dipanggil dari
keadaan nyawa. Memindahkan panggilan bunyi ke \code{FixedUpdate} menunda bunyi
satu sampai tiga bingkai, sebab bingkai fisika berjalan lebih jarang daripada
bingkai tampilan.

\subsection*{Kekuatan bunyi disusun, bukan dinaikkan satu per satu}

Kekuatan pada \istilah{AudioMixer} dinyatakan dalam desibel, dan skalanya tidak
lurus. Menurunkan 6 dB berarti mengurangi kekuatan menjadi kira-kira separuh.
Karena itu, menyeimbangkan bunyi lebih sering berarti menurunkan musik daripada
menaikkan efek suara. Selisih sekitar 10 dB antara efek dan musik adalah titik
awal yang baik.

% ---------------------------------------------------------
\bagian{Tutorial}

Nama menu dan jendela ditulis seperti yang tampak pada Unity LTS. Nama berkas
bunyi pada sesi ini hanya contoh; sesuaikan dengan berkas yang Anda miliki.

\begin{langkah}
  \item \textbf{Kumpulkan berkasnya.} Sediakan satu musik latar berdurasi
        panjang dan tiga efek suara pendek: lompat, mengambil benda, dan
        kalah. Catat nama berkas, sumber, dan lisensinya pada Lembar Kerja.

  \item \textbf{Buat folder dan impor.} Pada jendela \menu{Project}, klik
        kanan folder \file{Assets}, pilih \menu{Create > Folder}, dan beri
        nama \code{Audio}. Salin berkas bunyi ke folder itu melalui
        \menu{File Explorer}; Unity mengimpornya otomatis.

  \item \textbf{Atur impor efek suara.} Pilih berkas \code{lompat} pada
        jendela \menu{Project}. Pada \menu{Inspector}, setel
        \menu{Load Type} menjadi \menu{Decompress On Load},
        \menu{Compression Format} menjadi \menu{PCM}, dan centang
        \menu{Preload Audio Data}, lalu tekan \menu{Apply}. Ulangi untuk
        \code{ambil} dan \code{kalah}. Pengaturan ini menghilangkan penundaan
        ketika bunyi pendek dipanggil.

  \item \textbf{Atur impor musik.} Pilih berkas \code{musik-gudang}. Setel
        \menu{Load Type} menjadi \menu{Streaming} dan
        \menu{Compression Format} menjadi \menu{Vorbis}, lalu tekan
        \menu{Apply}. Musik berdurasi panjang tidak perlu dimuat seluruhnya ke
        memori.

  \item \textbf{Buat pemusat bunyi.} Klik kanan pada \menu{Hierarchy}, pilih
        \menu{Create Empty}, beri nama \code{ManajerBunyi}, dan letakkan di
        akar scene, bukan di bawah \code{Gudang}. Pilih
        \menu{Add Component} dua kali dan tambahkan \menu{Audio Source} dua
        kali.

  \item \textbf{Atur kedua sumber.} Pada \code{AudioSource} pertama, matikan
        \menu{Play On Awake}, centang \menu{Loop}, dan setel
        \menu{Spatial Blend} menjadi 0. Pada \code{AudioSource} kedua, matikan
        \menu{Play On Awake} dan \menu{Loop}, lalu setel
        \menu{Spatial Blend} menjadi 0 juga. Keduanya datar sebab bunyi
        permainan ini selalu berkaitan dengan pemain.

  \item \textbf{Buat AudioMixer.} Pada jendela \menu{Project}, klik kanan
        folder \file{Assets/Audio}, pilih \menu{Create > Audio Mixer}, dan
        beri nama \code{GudangMixer}. Klik dua kali berkas itu untuk membuka
        \menu{Window > Audio > Audio Mixer}. Pada panel \menu{Groups}, klik
        kanan \code{Master}, pilih \menu{Add Child Group} dua kali, lalu beri
        nama \code{Musik} dan \code{Efek}.

  \item \textbf{Atur tingkat dan keluaran.} Pilih tingkat \code{Musik}, setel
        \menu{Volume} menjadi $-8$ dB. Pilih tingkat \code{Efek}, biarkan
        pada 0 dB. Kembali ke \code{ManajerBunyi}, arahkan \menu{Output} pada
        \code{AudioSource} pertama ke \code{Master/Musik} dan \code{AudioSource}
        kedua ke \code{Master/Efek}.

  \item \textbf{Buat skrip pemusat.} Klik kanan folder \file{Assets/Scripts},
        pilih \menu{Create > C\# Script}, beri nama \code{ManajerBunyi}, lalu
        salin isi pada bagian Kode Sesi Ini. Seret skrip itu ke objek
        \code{ManajerBunyi}.

  \item \textbf{Hubungkan rujukannya.} Pada \menu{Inspector} komponen
        \code{ManajerBunyi}, isi \menu{Music Source} dan \menu{Sfx Source}
        dengan kedua \code{AudioSource} tadi. Isi \menu{Musik Latar},
        \menu{Bunyi Lompat}, \menu{Bunyi Ambil}, dan \menu{Bunyi Kalah} dengan
        klip yang sesuai. Setel \menu{Volume Musik} 0.4, \menu{Volume Efek}
        0.8, dan \menu{Variasi Nada} 0.08.

  \item \textbf{Buat penghubung pada pemain.} Buat skrip baru bernama
        \code{PemicuBunyi} dan salin isinya dari bagian Kode Sesi Ini. Seret
        ke objek \code{Pemain}, lalu isi \menu{Manajer} dengan objek
        \code{ManajerBunyi}.

  \item \textbf{Perbarui tiga skrip.} Ganti isi \file{KontrolPemain.cs},
        \file{Artefak.cs}, dan \file{TampilanHud.cs} dengan versi pada bagian
        Kode Sesi Ini. Pada \code{Pemain}, isi \menu{Pemicu Bunyi} dengan
        komponen \code{PemicuBunyi}. Pada setiap artefak dan pada
        \code{Canvas}, isi kolom \menu{Pemicu Bunyi} dengan komponen yang
        sama.
\end{langkah}

\begin{ujibox}
Setelah Langkah 12, periksa lima hal berikut.
\begin{itemize}[leftmargin=1.4em, itemsep=0.2em]
  \item Musik latar terdengar sejak scene dijalankan dan berulang setelah
        selesai, tanpa terputus ketika pemain bergerak.
  \item Bunyi lompat terdengar seketika ketika \menu{Space} ditekan, bersamaan
        dengan animasi lompat.
  \item Setiap artefak yang diambil mengeluarkan bunyi, dan dua artefak yang
        diambil berdekatan menghasilkan dua bunyi yang tidak saling memotong.
  \item Ketika nyawa mencapai nol, musik berhenti dan bunyi kalah terdengar.
  \item Tidak ada pesan merah atau peringatan rujukan kosong pada
        \menu{Console}.
\end{itemize}
\end{ujibox}

\begin{langkah}
  \item \textbf{Uji waktu dan kekuatan.} Jalankan scene, lalu ubah satu nilai
        setiap kali. Pertama, setel \menu{Volume Musik} 0.15 lalu 0.7 dan
        catat kapan efek suara mulai tertelan. Kedua, setel
        \menu{Variasi Nada} 0 lalu 0.25 dan bandingkan bunyi lompat yang
        dipanggil berulang. Ketiga, pindahkan sementara panggilan bunyi lompat
        ke \code{FixedUpdate} dan bandingkan keterlambatannya. Kembalikan
        setiap nilai setelah dicatat.

  \item \textbf{Simpan scene.} Tekan \menu{Ctrl+S}, lalu periksa kembali
        seluruh mekanik sesi 3: berjalan, melompat, mengambil artefak, dan
        membuka pintu.
\end{langkah}

\begin{ujibox}
Setelah Langkah 14, periksa dua hal berikut.
\begin{itemize}[leftmargin=1.4em, itemsep=0.2em]
  \item Pada \menu{Volume Musik} 0.15, efek suara terdengar jelas; pada 0.7,
        bunyi lompat mulai tertelan musik.
  \item Pada \menu{Variasi Nada} 0, bunyi yang dipanggil berulang terdengar
        seperti mesin; pada 0.25, nadanya berubah-ubah.
\end{itemize}
\end{ujibox}

% ---------------------------------------------------------
\bagian{Kode Sesi Ini}

\code{ManajerBunyi} menyimpan kedua sumber, seluruh klip, dan tingkat kekuatan,
lalu menyediakan satu fungsi untuk setiap peristiwa. \code{PemicuBunyi}
menghubungkan skrip lain dengan pemusat itu tanpa harus mengetahui isinya.
\code{KontrolPemain} memanggil bunyi lompat pada peristiwa masukan.
\code{Pintu} dan \code{AnimasiPemain} tidak berubah; salinannya tetap
disertakan pada checkpoint.

\tampilkode[language=CSharp,
  basicstyle=\ttfamily\scriptsize\color{ubCodeText}]
  {code/sesi-04/Scripts/ManajerBunyi.cs}
  {ManajerBunyi.cs: pemusat musik latar dan efek suara}

\tampilkode[language=CSharp,
  basicstyle=\ttfamily\scriptsize\color{ubCodeText}]
  {code/sesi-04/Scripts/PemicuBunyi.cs}
  {PemicuBunyi.cs: penghubung peristiwa permainan dengan bunyinya}

\tampilkodebaris[CSharp]
  {code/sesi-04/Scripts/KontrolPemain.cs}{10}{52}
  {KontrolPemain.cs: bidang bunyi dan Update yang memanggil bunyi lompat;
   bagian FixedUpdate tidak berubah dari sesi 3}

Dua skrip lain hanya berubah sedikit, sehingga isinya tidak dicetak ulang.
Pada \code{Artefak}, tambahkan dua peubah di bawah \code{hud}:

\begin{quote}
\code{public PemicuBunyi pemicuBunyi;}\\
\code{public ParticleSystem partikelAmbil;}
\end{quote}

lalu panggil keduanya di dalam \code{OnTriggerEnter2D} sebelum artefak
dimatikan:

\begin{quote}
\code{if (pemicuBunyi != null) pemicuBunyi.Ambil();}\\
\code{if (partikelAmbil != null) partikelAmbil.Play();}
\end{quote}

Medan partikel boleh dibiarkan kosong; bunyi tetap bekerja tanpa partikel.
Bila Anda menambahkannya, buat \menu{Effects > Particle System} terpisah di
posisi artefak, matikan \menu{Looping} dan \menu{Play On Awake}, lalu isi
medan itu.

Pada \code{TampilanHud}, tambahkan peubah berikut:

\begin{quote}
\code{public PemicuBunyi pemicuBunyi;}
\end{quote}

lalu panggil \code{pemicuBunyi.Kalah()} di dalam \code{KurangiNyawa} ketika
\code{lives} mencapai nol. Kedua berkas lengkap tersedia pada folder
checkpoint.


\subsection*{Tabel objek scene}

Tabel berikut adalah daftar objek yang harus ada pada scene Anda. Nilai yang
tidak disebut boleh dibiarkan bawaan.

{\footnotesize
\renewcommand{\arraystretch}{1.2}
\begin{tabularx}{\linewidth}{@{}L{2.4cm} L{3.4cm} Y@{}}
\toprule
\textbf{GameObject} & \textbf{Komponen} & \textbf{Nilai penting} \\
\midrule
ManajerBunyi & ManajerBunyi, dua AudioSource &
  Music Source dan Sfx Source terisi; empat klip terisi; Volume Musik 0.4;
  Volume Efek 0.8; Variasi Nada 0.08. Sumber musik: Loop ya, Spatial Blend 0,
  Output \code{Master/Musik}. Sumber efek: Loop tidak, Spatial Blend 0,
  Output \code{Master/Efek} \\
Pemain & KontrolPemain, PemicuBunyi &
  Pemicu Bunyi: komponen \code{PemicuBunyi}; Manajer: \code{ManajerBunyi} \\
Canvas & TampilanHud &
  Pemicu Bunyi: komponen \code{PemicuBunyi} pada pemain \\
Artefak1--3 & Artefak &
  Pemicu Bunyi terisi; Partikel Ambil boleh kosong \\
Pintu & Pintu &
  Tidak berubah dari sesi 3; Artefak Diperlukan 3 \\
\bottomrule
\end{tabularx}
}

\subsection*{Tabel peristiwa dan bunyi}

{\footnotesize
\renewcommand{\arraystretch}{1.2}
\begin{tabularx}{\linewidth}{@{}L{2.9cm} L{2.1cm} C{1.0cm} L{2.4cm} Y@{}}
\toprule
\textbf{Peristiwa} & \textbf{Klip} & \textbf{Jenis} & \textbf{Dipanggil oleh} &
\textbf{Pasangan yang terlihat} \\
\midrule
Scene dimulai & \code{musik-gudang} & BGM &
  \code{ManajerBunyi} & Level langsung dapat dimainkan \\
Tombol lompat & \code{lompat} & SFX &
  \code{KontrolPemain} & Animasi lompat dan posisi pemain naik \\
Artefak tersentuh & \code{ambil} & SFX &
  \code{Artefak} & Skor bertambah dan artefak lenyap \\
Nyawa mencapai nol & \code{kalah} & SFX &
  \code{TampilanHud} & Teks nyawa 0 dan ikon nyawa lenyap \\
\bottomrule
\end{tabularx}
}

\checkpoint{04}{\raggedright Folder checkpoint memuat
\file{Scripts/KontrolPemain.cs}, \file{Scripts/Artefak.cs},
\file{Scripts/Pintu.cs}, \file{Scripts/AnimasiPemain.cs},
\file{Scripts/TampilanHud.cs}, \file{Scripts/ManajerBunyi.cs},
\file{Scripts/PemicuBunyi.cs}, dan \file{scene.md}. Scene dan berkas bunyi
tidak disertakan sebagai berkas biner; susunlah mengikuti tabel di atas.}

% ---------------------------------------------------------
\bagian{Kesalahan Umum}

\begin{galatbox}
\begin{itemize}[leftmargin=1.4em, itemsep=0.35em]
  \item \textbf{Tidak ada bunyi sama sekali.}
        Gejala: scene berjalan tanpa suara, Console bersih. Periksa
        \menu{Volume} pada \code{AudioSource} dan pada tingkat
        \istilah{AudioMixer} (jangan sampai $-80$ dB), \menu{Mute} pada
        \code{AudioListener} kamera, serta klip yang belum diisi. Periksa juga
        \menu{Output} yang menunjuk tingkat senyap.
  \item \textbf{Musik dan efek suara saling memotong.}
        Gejala: musik terputus ketika pemain melompat, atau dua efek suara
        berdekatan saling memotong. Periksa jumlah \code{AudioSource}: musik
        dan efek suara harus memakai sumber berbeda. Periksa pula
        pemanggilannya: gunakan \code{PlayOneShot} untuk efek suara, bukan
        \code{Play}, dan beri \menu{Variasi Nada} nilai bukan nol.
  \item \textbf{Bunyi lompat berbunyi berulang selama tombol ditahan.}
        Gejala: bunyi berbunyi setiap bingkai. Periksa bahwa panggilan berada
        di dalam \code{Input.GetButtonDown}, bukan di luar blok itu atau pada
        \code{FixedUpdate} berdasarkan keadaan tombol.
  \item \textbf{Bunyi terdengar terlambat.}
        Gejala: bunyi muncul beberapa saat setelah lompatan. Periksa tempat
        pemanggilan: bunyi harus dipanggil di \code{Update} pada peristiwa
        masukan, bukan di \code{FixedUpdate}. Periksa pula
        \menu{Load Type} klip pendek; \menu{Streaming} pada efek suara
        menambah penundaan.
  \item \textbf{Bunyi kalah tidak pernah terdengar.}
        Gejala: nyawa habis tetapi musik tetap berjalan. Periksa rujukan
        \menu{Pemicu Bunyi} pada \code{Canvas}, lalu uji melalui menu konteks
        \menu{Uji kurangi satu nyawa} pada komponen \code{TampilanHud}.
        Bila Console menampilkan \code{NullReferenceException}, periksa pula
        rujukan \menu{Manajer} pada \code{PemicuBunyi} dan \menu{Music Source}
        pada \code{ManajerBunyi}.
\end{itemize}
\end{galatbox}

% ---------------------------------------------------------
\bagian{Lembar Kerja}

Isi tabel berikut, lalu simpan sebagai bagian dari catatan praktikum Anda.

{\small
\renewcommand{\arraystretch}{1.4}
\begin{tabularx}{\linewidth}{@{}C{0.8cm} L{6.2cm} Y@{}}
\toprule
\textbf{No} & \textbf{Hasil yang dicatat} & \textbf{Keterangan Anda} \\
\midrule
1 & Nama, sumber, dan lisensi setiap berkas bunyi & \\[0.5cm]
2 & Pengaturan impor setiap berkas (Load Type, Compression) & \\[0.5cm]
3 & Nilai Volume Musik, Volume Efek, dan Variasi Nada yang dipakai & \\[0.5cm]
4 & Volume tingkat Musik dan Efek pada AudioMixer & \\[0.5cm]
5 & Hasil percobaan Volume Musik 0.15 dan 0.7 & \\[0.9cm]
6 & Hasil percobaan Variasi Nada 0 dan 0.25 & \\[0.9cm]
7 & Satu masalah bunyi yang muncul dan cara memperbaikinya & \\[0.9cm]
\bottomrule
\end{tabularx}
}

% ---------------------------------------------------------
\bagian{Checklist Sesi}

\begin{checklistbox}
\begin{checklist}
  \item Saya dapat menunjukkan \istilah{AudioClip} yang terimpor dengan
        pengaturan \menu{Load Type} yang sesuai untuk musik dan efek suara.
  \item Saya dapat menjelaskan perbedaan \istilah{AudioClip} dan
        \istilah{AudioSource}, serta alasan keduanya dipisahkan.
  \item Saya dapat menunjukkan dua \istilah{AudioSource} terpisah untuk
        \istilah{BGM} dan \istilah{SFX}, beserta keluaran masing-masing pada
        \istilah{AudioMixer}.
  \item Saya dapat menunjukkan bunyi lompat, bunyi mengambil artefak, dan bunyi
        kalah yang terdengar pada peristiwanya.
  \item Saya dapat menjelaskan mengapa bunyi dipanggil di \code{Update} dan
        bukan di \code{FixedUpdate}.
  \item Saya dapat menyeimbangkan musik dan efek suara dengan mengubah tingkat
        pada \istilah{AudioMixer}, bukan dengan menaikkan setiap berkas.
\end{checklist}
\end{checklistbox}

% ---------------------------------------------------------
\bagian{Pengait ke Sesi Berikutnya}

Level kini terlihat dan terdengar lengkap. Yang belum dikerjakan adalah
memeriksa apakah seluruh bagian itu benar-benar tersambung, dan apakah pemain
lain dapat menyelesaikannya tanpa penjelasan Anda. Sesi 5 menuntun penyusunan
\istilah{checklist} fungsional, pelaksanaan \istilah{playtest} antar kelompok,
pencatatan \istilah{log bug}, dan presentasi demo. Di sanalah bunyi yang terlalu
pelan atau datang terlambat akan ditemukan oleh orang lain, bukan oleh Anda
sendiri.
